    \documentclass[conference]{IEEEtran}
    \usepackage{cite}
    \usepackage{amsmath,amssymb,amsfonts}
    \usepackage{algorithmic}
    \usepackage{graphicx}
    \usepackage{textcomp}
    \usepackage{xcolor}
    \usepackage{booktabs}
    \usepackage{url}

    \begin{document}

    \title{Tether-Edge: AI Orchestration for Distributed Edge Devices}

    \author{
    \IEEEauthorblockN{Anish Chakrabarti, Joshua DeNiese, Mohammad Sufiyan}
    \IEEEauthorblockA{\textit{Department of Computer Science \& Engineering} \\
    \textit{MIT-ADT University, MIT School of Computing}\\
    Pune, Maharashtra, India \\}
    \and
    \IEEEauthorblockN{Prof. Dr. Harshad Lokhande}
    \IEEEauthorblockA{\textit{Faculty Guide} \\
    \textit{MIT-ADT University}\\
    Pune, Maharashtra, India}
    }

    \maketitle

    \begin{abstract}
    AI is increasingly being installed directly onto physical hardware---such as security cameras, factory machinery, and robots---so they can process sensor data
    on the spot without relying on the cloud. However, managing hundreds or thousands of these devices is very challenging. These devices have limited memory and
    battery power, frequently lose internet connectivity, and traditional software management tools are often too heavy and slow. On many systems, everything runs
    as a single large program; if an AI model crashes or runs out of memory, the entire device goes down and critical data is lost. Furthermore, running heavy AI
    models continuously causes devices to overheat and quickly drains their batteries.

    To solve these problems, this paper introduces Tether-Edge, a lightweight and reliable management platform designed specifically for fleets of edge AI devices.
    Instead of running everything as one big program, Tether-Edge splits tasks into small, separate, and safe components that communicate with each other smoothly.
    If an AI workload or rule crashes, it is safely contained so the device stays online and continues operating normally.

    To save battery and prevent overheating, Tether-Edge uses a fast, lightweight filter that quickly checks incoming sensor data. It ignores normal, repetitive
    scenes and only wakes up the heavy AI models when an actual change or physical problem is detected. Once triggered, the system can spot objects on camera,
    check for abnormal machine vibrations or unusual sounds, and even generate a plain-English explanation of what went wrong using a built-in, on-device AI
    assistant---all without sending data to an external cloud.

    To handle unstable internet connections, Tether-Edge safely stores all reports and sensor logs on the device itself and uploads them automatically once the
    network reconnects, ensuring zero data loss. In addition, when new AI models are updated remotely over the air, the system smoothly swaps in the new version in
    the background without restarting the device or dropping a single camera frame. Tests show that Tether-Edge requires very little memory, cuts unnecessary
    processing work, and runs reliably on small, low-cost edge computers.
    \end{abstract}

    \begin{IEEEkeywords}
    Edge AI, Device Fleet Management, System Orchestration, Crash Protection, Smart Data Filtering, Multi-Sensor Inspection, Offline Reliability, Zero-Downtime
    Updates.
    \end{IEEEkeywords}

    \section{Introduction}
    Edge Artificial Intelligence (Edge AI) involves running machine learning models directly on physical edge devices---such as smart cameras, industrial 
    machinery, and robotic systems---rather than sending raw data to distant cloud data centers. By processing sensor streams locally, edge computing significantly
    reduces network latency, minimizes cellular data consumption, and preserves strict data privacy \cite{b1, b4}.

    However, deploying, updating, and operating large fleets of geographically distributed edge devices presents several critical challenges:
    \begin{enumerate}
        \item \textbf{Device Outages from Single-Component Failures}: Many edge systems operate as monolithic applications. If a third-party AI model encounters an
        unexpected software bug, memory leak, or crash, the entire application fails. This severs remote management connections and often requires an engineer to
        travel on-site simply to restart the device.
        \item \textbf{Excessive Infrastructure Overhead}: Mainstream container platforms, such as standard Docker or Kubernetes distributions, were primarily
        engineered for high-performance cloud servers. Running these frameworks consumes substantial memory (often 300\,MB to over 1\,GB) and processing resources,
        leaving insufficient headroom for AI workloads on low-power, entry-level hardware (e.g., 1--4\,GB RAM single-board computers) \cite{b6, b7}.
        \item \textbf{Overheating and Battery Depletion}: Executing continuous neural network inference across every incoming video frame or high-frequency sensor
        reading demands significant computational power. On fanless industrial enclosures or battery-operated platforms, this leads to rapid battery drain and 
        severe thermal throttling \cite{b5}.
        \item \textbf{Data Loss from Unstable Connectivity}: Remote edge hardware frequently connects through spotty cellular or satellite links. When
        cloud-dependent systems experience network drops, real-time alerts, error logs, and sensor telemetry are discarded and permanently lost.
        \item \textbf{Downtime and Missed Events During Updates}: In round-the-clock industrial environments, software updates typically require restarting
        services. This introduces several seconds of operational downtime, during which cameras stop analyzing data and critical defects or safety violations can
        pass unnoticed.
        \item \textbf{Bulky and Non-Deterministic Runtimes}: Early prototype implementations frequently depend on heavy interpreted environments (such as Python
        and PyTorch). These runtimes require gigabytes of disk storage and introduce unpredictable execution pauses, making them unsuitable for consistent, 
        real-time edge operations.
    \end{enumerate}

    To address these challenges, this paper presents \textbf{Tether-Edge}, a lightweight, reliable, and fault-tolerant orchestration platform tailored for
    distributed Edge AI fleets. Tether-Edge replaces monolithic architectures with small, decoupled, and memory-safe native Rust microservices that communicate
    over an ultra-fast local message bus. To conserve energy and prevent thermal throttling, the platform uses a lightweight WebAssembly sandbox to pre-filter
    sensory data, only triggering deep neural models when physical state changes are detected. Furthermore, Tether-Edge incorporates an offline store-and-forward
    telemetry buffer to prevent data loss during network disconnections, a built-in on-device language model for local root-cause diagnostics, and a zero-downtime
    model swap mechanism that applies updates without dropping a single sensor frame.

    \section{Literature Survey}

    \subsection{Edge AI and Edge Computing}
    Seminal research by Shi et al. \cite{b1} and Satyanarayanan \cite{b4} demonstrated that processing data locally on edge hardware resolves the latency and
    bandwidth limitations of centralized cloud architectures. Subsequent work by Bonomi et al. \cite{b2} and the OpenFog Consortium \cite{b3} standardized
    distributed fog computing frameworks. However, running deep neural networks continuously on low-power hardware leads to rapid battery exhaustion and thermal
    throttling, highlighting the need for efficient pre-filtering mechanisms that activate heavy AI models only when necessary.

    \subsection{Containerization vs. WebAssembly Isolation}
    Docker containerization popularized microservice isolation by bundling application dependencies into portable images \cite{b7, b8}. However, empirical studies
    by Paul et al. \cite{b13} revealed that running container orchestrators (such as K3s or MicroK8s) on resource-constrained hardware can consume up to 40\% of
    available system memory simply maintaining the orchestration layer. In response, WebAssembly (WASM) and the WebAssembly System Interface (WASI) have emerged as
    lightweight sandboxing alternatives \cite{b12}. WebAssembly runtimes launch within milliseconds and enforce strict memory limits with minimal resource
    overhead. Tether-Edge uses compiled WebAssembly modules to execute user rules and filter sensor data without the high memory footprint of traditional
    containers.

    \subsection{Edge Orchestration and Fleet Management}
    Commercial platforms such as AWS IoT Greengrass \cite{b9} and Azure IoT Edge \cite{b10} offer device provisioning and remote application deployment. However,
    these systems introduce proprietary vendor lock-in and require heavy local runtimes. Earlier distributed systems like Condor \cite{b11} provided resilient
    remote job execution, but targeted large, high-performance computing clusters rather than resource-limited edge nodes. Tether-Edge fills this gap by providing
    an open, lightweight architecture built on fast local messaging and local database buffering.

    \subsection{Inter-Process Communication and Message Buses}
    Microservices frequently communicate using local HTTP or network sockets, which introduce processing delays and operating system context-switching overhead.
    Utilizing a lightweight publish/subscribe message broker (such as NATS over local system sockets) significantly reduces communication latency, decouples
    individual software services, and ensures that an isolated application crash does not destabilize the broader system.

    \subsection{Multi-Modal Industrial Sensing}
    Effective industrial monitoring relies on combining complementary sensor inputs. Vision-based models (such as YOLO) provide real-time object and surface defect
    detection, while vibration analysis conforming to international standards (ISO-10816) can identify mechanical imbalance and bearing degradation \cite{b5}. In
    addition, acoustic analysis can detect early signs of pump friction and wear. Tether-Edge synchronizes visual, vibrational, and acoustic data into a unified
    edge data stream for comprehensive equipment monitoring.

    \subsection{On-Device Small Language Models (SLMs)}
    While Large Language Models have shown outstanding diagnostic capabilities, serving them historically required powerful data center GPUs \cite{b14}. Recent
    advancements in model quantization and lightweight inference software \cite{b15} allow compact Small Language Models to execute efficiently on standard
    low-power processors. Tether-Edge incorporates an on-device language model to generate clear, natural language root-cause explanations directly at the edge
    without requiring external network connectivity.

    \section{Objectives and Scope}

    \subsection{Objectives}
    The primary objective of Tether-Edge is to design, implement, and validate an industrial-grade, lightweight fleet orchestration platform for distributed Edge
    AI systems. The specific research and engineering objectives are:
    \begin{itemize}
        \item \textbf{Minimal Memory Footprint}: Maintain the total steady-state memory utilization of all edge client services below \textbf{520\,MB RAM},
        enabling deployment on cost-effective edge devices equipped with 1--4\,GB of RAM.
        \item \textbf{Rapid Sensor Event Gating}: Implement an ultra-fast WebAssembly (WASM) pre-filtering sandbox that inspects incoming data in under 0.02\,ms,
        filtering out over 90\% of nominal sensor frames to reduce downstream AI processing load by more than 70\%.
        \item \textbf{Pure Native Edge Execution}: Eliminate heavy interpreted runtime environments (such as Python and virtual environments) on edge devices,
        compiling the entire client suite into lightweight, memory-safe native Rust and C++ binaries.
        \item \textbf{Synchronized Multi-Modal Ingestion}: Simultaneously capture and analyze live optical camera streams, multi-axis machine vibration signals,
        acoustic audio, and system thermal metrics within a synchronized pipeline.
        \item \textbf{Zero Data Loss via Offline Buffering}: Implement a local store-and-forward SQLite database that safely buffers telemetry and anomaly alerts
        during network outages, automatically synchronizing data once connectivity is restored.
        \item \textbf{On-Device Diagnostic Reasoning}: Integrate a compact, 4-bit quantized language model capable of diagnosing equipment faults and generating
        real-time operator guidance directly on the local processor without external network access.
        \item \textbf{Zero-Downtime Over-The-Air (OTA) Updates}: Enable dynamic AI model updates in the background without requiring service restarts, ensuring 
        continuous sensor ingestion and zero dropped camera frames during model transitions.
    \end{itemize}

    \subsection{Scope and Feasibility}
    The scope of Tether-Edge encompasses resource-constrained edge hardware equipped with 1--4\,GB of RAM, including single-board computers (e.g., Raspberry Pi
    4/5), embedded edge platforms (e.g., NVIDIA Jetson), and standard industrial x86 gateways. The platform manages device registration, multi-sensor data capture,
    sandboxed rule execution, dynamic AI model updates, offline telemetry storage, and real-time fleet dashboard monitoring.

    The project demonstrates high technical feasibility by leveraging mature, proven open-source technologies rather than requiring proprietary cloud runtimes or
    specialized hardware security chips:
    \begin{itemize}
        \item \textbf{Rust}: Delivers native execution speed and strict memory safety without the overhead of runtime garbage collection.
        \item \textbf{Wasmtime}: Provides lightweight, capability-based WebAssembly sandboxing to execute user rules safely with minimal memory usage.
        \item \textbf{Microsoft ONNX Runtime}: Enables optimized, real-time neural network inference directly on low-power central processing units.
        \item \textbf{SQLite}: Offers reliable local database buffering in Write-Ahead Logging mode to ensure persistent telemetry storage during network drops.
        \item \textbf{NATS}: Supplies a fast, asynchronous publish/subscribe message broker for decoupled inter-process communication.
        \item \textbf{\texttt{llama.cpp}}: Facilitates efficient on-device execution of quantized language models for real-time diagnostic reasoning without
        internet access.
    \end{itemize}

    \section{Proposed System Architecture}

    \subsection{Architectural Overview}
    Tether-Edge acts as an intelligent intermediary between centralized cloud management and distributed edge hardware. The system follows a decoupled,
    event-driven microservices architecture partitioned into two primary operational zones: the \textbf{Edge Hardware Platform} and the \textbf{Cloud Control
    Plane}, as illustrated in Fig.~\ref{fig:architecture}.

    \begin{figure*}[!t]
    \centering
    \includegraphics[width=0.92\textwidth]{architecture_diagram.png}
    \caption{System architecture of the Tether-Edge platform, depicting the hardware abstraction layer, local message bus, WebAssembly pre-filter sandbox, native
    multi-modal AI engine, on-device language model, offline storage buffer, and cloud control plane.}
    \label{fig:architecture}
    \end{figure*}

    \subsection{Edge Microservices Breakdown}
    The edge client suite is divided into independent, unprivileged processes communicating exclusively over an ultra-fast local NATS message bus. This design
    guarantees fault isolation: if a specific worker fails or restarts, the remaining edge microservices continue operating without disruption.

    \subsubsection{\texttt{client\_daemon} (Master Orchestrator \& Hardware Layer)}
    Implemented in native Rust, this daemon supervises device operations and manages hardware inputs:
    \begin{itemize}
        \item \textbf{Optical Video}: Captures live video frames directly from camera hardware devices into shared memory and computes frame luminance variations.
        \item \textbf{System Thermals}: Monitors hardware thermal sensors via the Linux kernel to track operating temperatures.
        \item \textbf{Vibration and Audio}: Ingests multi-axis vibration signals and acoustic audio samples at native sampling rates.
        \item \textbf{Telemetry Dispatch}: Formats sensor payloads and publishes them to dedicated internal message topics while streaming real-time status to the
        cloud gateway.
    \end{itemize}

    \subsubsection{\texttt{wasm\_host} (High-Speed WASM Pre-Filtering Sandbox)}
    Powered by Wasmtime, this worker executes safety rules and pre-filters incoming sensor data in under 0.02\,ms. Custom rules are compiled into WebAssembly
    modules and executed within a strictly isolated memory sandbox. Host and guest exchange data via an efficient in-memory buffer, eliminating unnecessary
    serialization overhead. Unchanged scenes and normal sensor readings are discarded immediately, preventing unnecessary downstream neural computation and saving
    processor power.

    \subsubsection{\texttt{inference\_worker} (Native Multi-Modal AI Engine)}
    Built entirely in native Rust, this service replaces heavy Python machine learning stacks with an optimized inference engine using the Microsoft ONNX Runtime
    (\texttt{ort}) library. To minimize system memory overhead on resource-constrained devices, Tether-Edge applies five key architectural optimizations:
    \begin{enumerate}
        \item \textbf{Thread Pool Clamping}: Limits internal execution threads to prevent excessive thread stack memory allocation and reduce operating system
        context switching.
        \item \textbf{Deterministic Execution}: Replaces dynamic parallel scheduling with sequential operator execution, maintaining predictable memory
        utilization.
        \item \textbf{Static Memory Reuse}: Pre-allocates intermediate memory buffers based on the model structure, reusing memory across consecutive layers and
        preventing heap fragmentation.
        \item \textbf{Graph Optimization}: Fuses consecutive neural network layers during model initialization to accelerate processing speed and streamline
        execution.
        \item \textbf{Zero-Copy Tensor Views}: Directly references existing image buffers in memory rather than duplicating tensor data across the system heap.
    \end{enumerate}
    In addition to computer vision, this worker handles vibration analysis, acoustic anomaly detection, and automated queries to the local diagnostic language
    model.

    \subsubsection{\texttt{storage\_service} (Offline Buffer \& Staging Gatekeeper)}
    Implemented in Rust using SQLite in Write-Ahead Logging (WAL) mode, this service provides high-speed local data logging without lock contention, while also
    managing safe model staging during remote updates.

    \subsubsection{\texttt{llama-server} (On-Device Language Model)}
    A lightweight, native executable based on \texttt{llama.cpp} that loads 4-bit quantized language models (e.g., Qwen-0.6B). It generates natural language 
    diagnostic reports and operator guidance locally on the device CPU without requiring external internet access.

    \subsubsection{\texttt{cloud\_server} (Control Plane \& Fleet Dashboard)}
    A lightweight web server providing device registration, secure authentication, remote model cataloging, and an interactive fleet dashboard featuring a visual
    pipeline canvas and live video inspector with bounding box overlays.

    \subsection{Inter-Process Communication (IPC)}
    Monolithic edge systems frequently crash entirely when a single AI model encounters a fault. Tether-Edge prevents this fate-sharing vulnerability by deploying
    an asynchronous local \textbf{NATS message bus}. Individual services operate as independent producers and consumers across dedicated channels:
    \begin{itemize}
        \item Raw sensor channels for camera frames, vibration waveforms, and audio signals.
        \item Filtered event and anomaly alert channels emitted by the WebAssembly sandbox.
        \item Inference request and diagnostic output channels.
        \item Fleet administrative channels for remote updates, health checks, and heartbeat monitoring.
    \end{itemize}

    \subsection{Offline SQLite Telemetry Buffer \& Store-and-Forward Sync}
    Direct cloud HTTP connections routinely drop packets during intermittent network outages on industrial networks. Tether-Edge implements a store-and-forward
    architecture managed by \texttt{storage\_service}. All sensor events, anomaly alerts, and diagnostic outputs are appended immediately to a local database file.
    During network disconnections, edge services continue operating and logging locally without blocking. Once connectivity is re-established, the system
    automatically synchronizes buffered records to the cloud, guaranteeing zero telemetry loss.

    \subsection{Cryptographic Zero-Downtime Drain-and-Swap OTA Protocol}
    Updating AI models on remote edge devices typically requires stopping containers, leading to service downtime and dropped video frames. Tether-Edge implements
    a \textbf{Drain-and-Swap} hot-reloading lifecycle with four distinct reliability guarantees:
    \begin{enumerate}
        \item \textbf{Staged Download}: The new model binary is downloaded to a temporary staging location while the active model continues executing.
        \item \textbf{Cryptographic Verification}: The file's SHA-256 checksum is validated against the cloud manifest; corrupted or incomplete downloads are
        rejected immediately.
        \item \textbf{Atomic File Commit}: The validated binary is committed using an atomic file rename operation, ensuring that the system never accesses a
        partially written file on disk.
        \item \textbf{Queue Draining and Pointer Swapping}: Upon receiving an update notification, the inference service temporarily pauses queue consumption while
        in-flight frames finish execution. The worker loads the new model into memory, updates the active pointer, and resumes processing immediately. This process
        completes with \textbf{zero dropped frames and zero service restarts}.
    \end{enumerate}

    \section{Methodology and Implementation Details}

    \subsection{Multi-Modal Sensing and Detection Logic}

    \subsubsection{Optical Frame Gating and Surface Defect Detection}
    To prevent energy-intensive neural networks from executing continuously on static scenes, the camera layer calculates temporal luminance changes across 
    consecutive video frames. If the mean pixel variation between frames remains below a 4\% motion threshold, the scene is classified as static. The WebAssembly
    pre-filter suppresses downstream processing, preventing unnecessary AI compute cycles. When motion is detected, the pipeline evaluates spatial contrast
    gradients across pixel boundaries; sudden high-contrast variations exceeding threshold limits are flagged as potential surface fissures or scratches.

    \subsubsection{Tri-Axial Vibration and ISO-10816 Severity Classification}
    Machine health is monitored using three-axis vibration velocity evaluated against the international predictive maintenance standard ISO-10816. The vibration
    engine calculates the Root Mean Square (RMS) velocity across orthogonal axes and applies a discrete high-frequency impact filter to identify repetitive shock
    pulses:
    \begin{itemize}
        \item \textbf{Class I (Nominal, $\le 1.8$\,mm/s)}: Smooth mechanical operation with no discernible bearing degradation.
        \item \textbf{Class II (Warning, $1.8\text{--}4.5$\,mm/s)}: Elevated vibration indicating shaft misalignment or rotor unbalance.
        \item \textbf{Class III/IV (Critical, $>4.5$\,mm/s)}: Excessive mechanical shock and high-frequency harmonic energy, indicating severe bearing degradation.
    \end{itemize}

    \subsubsection{Acoustic Cavitation and Friction Analysis}
    Fluid handling equipment is monitored through acoustic frequency analysis. The audio pipeline computes the energy ratio between ultrasonic high-frequency bands
    and baseline operational noise. A sustained energy shift toward high frequencies signals fluid turbulence, pump cavitation, or mechanical friction rub, which
    automatically triggers diagnostic analysis.
    
    \section{Performance Evaluation and Empirical Results}
    
    \subsection{Experimental Testbed}
    All empirical evaluations were conducted on a physical Linux edge node under realistic operating conditions, as detailed in Table~\ref{tab:testbed}.
    
    \begin{table}[htbp]
    \caption{Physical Edge Testbed Specifications}
    \label{tab:testbed}
    \centering
    \resizebox{\columnwidth}{!}{%
    \begin{tabular}{ll}
    \toprule
    \textbf{Parameter} & \textbf{Specification} \\
    \midrule
    Edge Processor & 11th Gen Intel Core i5-1135G7 (4 Cores / 8 Threads, up to 4.2\,GHz) \\
    System Memory & 8\,GB DDR4 (Software target budget $<520$\,MB) \\
    Camera Interface & Physical UVC Webcam (\texttt{/dev/video0}) \\
    Video Stream & 640$\times$480 @ 30\,FPS (Direct V4L2 MJPEG) \\
    Operating System & Ubuntu Linux 24.04 LTS (x86\_64) \\
    WASM Runtime & Bytecode Alliance Wasmtime 21.0 (\texttt{wasm32-wasip1}) \\
    Inference Engine & Pure Rust with Microsoft ONNX Runtime (\texttt{ort 2.0}) \\
    Local SLM & \texttt{llama-server} with Qwen-0.6B Q4\_K\_M GGUF \\
    IPC Broker & Local NATS Server v2.10.14 \\
    Storage Engine & SQLite 3.45 in Write-Ahead Logging (WAL) Mode \\
    \bottomrule
    \end{tabular}%
    }
    \end{table}
    
    \subsection{Tiered WASM Gating \& Computational Savings}
    We evaluated the efficiency of the WebAssembly pre-filter sandbox during a continuous optical monitoring benchmark simulating an industrial environment with
    intermittent activity (12\% active duty cycle).
    
    \begin{table}[htbp]
    \caption{System Performance: Continuous Inference vs. Tiered Gating}
    \label{tab:gating}
    \centering
    \resizebox{\columnwidth}{!}{%
    \begin{tabular}{lll}
    \toprule
    \textbf{Metric} & \textbf{Continuous Baseline} & \textbf{Tiered WASM Gating} \\
    \midrule
    Total Processed Frames & 108,000 frames & 108,000 frames \\
    Pre-Filter Gating Latency & None (0\,ms) & \textbf{$<$0.02\,ms (WASM Sandbox)} \\
    Neural Network Invocations & 108,000 passes & \textbf{12,960 passes (-88.0\%)} \\
    Static Frames Filtered & 0 frames & 95,040 frames (100\% accurate) \\
    Mean CPU Utilization & 78.4\% & \textbf{9.1\% (-88.4\% load)} \\
    Peak Package Temperature & 74.2\,$^{\circ}$C & \textbf{46.8\,$^{\circ}$C (No thermal throttling)} \\
    \bottomrule
    \end{tabular}%
    }
    \end{table}
    
    As shown in Table~\ref{tab:gating}, continuous inference caused high processor utilization (78.4\%) and elevated package temperatures (74.2\,$^{\circ}$C). In
    contrast, the WebAssembly pre-filter evaluated frames in under 0.02\,ms, filtering out 95,040 static frames and reducing active AI execution by 88.0\%, which
    dropped operating temperature by 27.4\,$^{\circ}$C and completely prevented thermal throttling.
    
    \subsection{Inference Engine Benchmark: Native Rust ORT vs. Python}
    We benchmarked the native Rust ONNX Runtime implementation against the traditional Python/PyTorch stack on identical 640$\times$480 camera frames.
    
    \begin{table}[htbp]
    \caption{Inference Engine Benchmark on CPU (640$\times$480 Frames)}
    \label{tab:inference}
    \centering
    \resizebox{\columnwidth}{!}{%
    \begin{tabular}{lll}
    \toprule
    \textbf{Metric} & \textbf{Python/PyTorch Bridge} & \textbf{Native Rust ORT} \\
    \midrule
    Runtime Environment & Python 3.12 + PyTorch & \textbf{Pure Rust (\texttt{ort 2.0})} \\
    Model Format & PyTorch \texttt{.pt} & \textbf{ONNX \texttt{.onnx}} \\
    Preprocessing Time & 18.4\,ms & \textbf{30.3\,ms} \\
    Tensor Forward Pass & 384.2\,ms & \textbf{124.3\,ms (SIMD Fused)} \\
    Post-Processing Time & 12.1\,ms & \textbf{2.4\,ms} \\
    Total Latency & 414.7\,ms & \textbf{157.0\,ms ($\sim$6.4\,FPS)} \\
    Resident Memory (RSS) & 428\,MB & \textbf{14.2\,MB (Zero Python)} \\
    External Dependencies & Python, pip, virtualenv & \textbf{None (Static ELF)} \\
    \bottomrule
    \end{tabular}%
    }
    \end{table}
    
    As shown in Table~\ref{tab:inference}, the native Rust ONNX engine completes a forward pass in 124.3\,ms compared to 384.2\,ms in the Python bridge. More
    importantly, eliminating the Python interpreter and virtual environment reduced the resident memory footprint from 428\,MB to just 14.2\,MB.
    
    \subsection{System-Wide Memory Footprint Audit}
    Table~\ref{tab:memory} presents the audited memory profile across both idle and peak active inference execution.
    
    \begin{table}[htbp]
    \caption{Comprehensive Edge Microservices Memory Audit (RSS)}
    \label{tab:memory}
    \centering
    \resizebox{\columnwidth}{!}{%
    \begin{tabular}{lllll}
    \toprule
    \textbf{Microservice Binary} & \textbf{Idle RSS} & \textbf{Peak RSS} & \textbf{Idle CPU} & \textbf{Active CPU} \\
    \midrule
    \texttt{client\_daemon} (HAL) & 13.9\,MB & 14.3\,MB & 1.2\% & 4.5\% \\
    \texttt{wasm\_host} (Wasmtime) & 8.2\,MB & 25.4\,MB & 0.2\% & 1.8\% \\
    \texttt{wasm\_worker.wasm} & 64\,KB & 64\,KB & 0.0\% & $<$0.02ms \\
    \texttt{inference\_worker} (ORT) & 14.2\,MB & \textbf{106.5\,MB} & 0.1\% & 28.4\% \\
    \texttt{storage\_service} (SQLite)& 3.4\,MB & 7.0\,MB & 0.1\% & 1.2\% \\
    \texttt{nats-server} (IPC Bus) & 11.5\,MB & 12.8\,MB & 0.3\% & 1.5\% \\
    \midrule
    \textbf{Core Orchestration Suite} & \textbf{51.2\,MB} & \textbf{166.0\,MB} & \textbf{1.9\%} & \textbf{37.4\%} \\
    \midrule
    \texttt{llama-server} (Qwen-0.6B)& 468.0\,MB & 859.4\,MB & 0.0\% & 98.2\% \\
    \midrule
    \textbf{Total Suite (with Local SLM)} & \textbf{519.2\,MB} & \textbf{1,025.4\,MB} & \textbf{1.9\%} & \textbf{135.6\%} \\
    \bottomrule
    \end{tabular}%
    }
    \end{table}
    
    By enforcing thread clamping and static memory reuse, the core orchestration suite executes within \textbf{166.0\,MB peak RAM}, leaving ample resources
    available on 1--4\,GB edge devices. Even when the local language model is actively generating diagnostic text, the entire system operates comfortably within
    1\,GB of memory.
    
    \subsection{Multi-Modal Diagnostic Latency}
    Table~\ref{tab:multimodal} summarizes the inference latency across all monitored modalities.
    
    \begin{table}[htbp]
    \caption{Multi-Modal Detection \& Diagnostic Latency Summary}
    \label{tab:multimodal}
    \centering
    \resizebox{\columnwidth}{!}{%
    \begin{tabular}{llll}
    \toprule
    \textbf{Modality} & \textbf{Target Metric} & \textbf{Backend Engine} & \textbf{Latency} \\
    \midrule
    Vision (Objects) & COCO Object Classes & \texttt{yolov8n.onnx} & 124.3\,ms \\
    Vision (Defects) & Surface Fissures & Spatial Convolution & 1.2\,ms \\
    Vibration & 3-Axis Bearing Wear & 1D-CNN (ISO-10816) & 1.5\,ms \\
    Acoustic & Pump Cavitation & Mel-Spectrum Ratio & 1.8\,ms \\
    On-Device SLM & Operator Diagnostics & Qwen-0.6B GGUF & 45.3\,tok/s \\
    \bottomrule
    \end{tabular}%
    }
    \end{table}
    
    When vibration exceeded the ISO-10816 critical threshold ($>4.5\text{ mm/s}$), the vibration engine classified bearing degradation in 1.5\,ms. The inference
    worker queried the local language model, which generated a 58-token plain-English diagnosis and repair recommendation in 1.28 seconds at
    \textbf{45.3\,tokens/second directly on CPU}, with zero external cloud dependencies.
    
    \section{Applications and Real-World Use Cases}
    The Tether-Edge platform addresses mission-critical applications across distributed environments:
    \begin{itemize}
        \item \textbf{Industrial Manufacturing \& Predictive Maintenance}: Orchestrates multi-modal inspection on factory floors, combining optical surface defect
        detection with three-axis vibration monitoring to prevent unexpected equipment downtime.
        \item \textbf{Smart Cities \& Intelligent Transportation}: Analyzes video streams directly on traffic cameras to monitor congestion and road hazards,
        minimizing cellular bandwidth usage by transmitting only essential event summaries.
        \item \textbf{Critical Infrastructure \& Energy Grids}: Monitors remote power substations, pipelines, and wind turbines where network coverage is
        intermittent, utilizing the on-device language model to deliver autonomous diagnostic guidance.
        \item \textbf{Healthcare \& Patient Monitoring}: Evaluates patient vitals and room cameras on local bedside gateways, ensuring sensitive health records
        never leave hospital premises.
        \item \textbf{Remote \& Tactical Field Deployments}: Maintains continuous situational awareness in disconnected environments where cloud access is
        unavailable.
    \end{itemize}
    
    \section{Expected Outcomes and Future Scope}
    
    \subsection{Expected Outcomes}
    Tether-Edge delivers a production-grade, lightweight fleet orchestration platform that ensures continuous edge device availability, eliminates runtime memory
    bloat, and prevents telemetry loss during network outages. System administrators can remotely deploy AI models and safety rules across distributed devices
    through a unified web dashboard, ensuring reliable performance, zero data loss, and zero-downtime updates.
    
    \subsection{Future Scope}
    Future enhancements for Tether-Edge include:
    \begin{enumerate}
        \item \textbf{Dedicated NPU Acceleration}: Integrating hardware Neural Processing Unit (NPU) backends to further accelerate inference on specialized mobile
        silicon.
        \item \textbf{Privacy-Preserving Federated Learning}: Enabling decentralized model training across edge fleets using local database storage without
        centralizing private data.
        \item \textbf{Deterministic Real-Time Scheduling}: Incorporating real-time Linux kernel scheduling extensions for mission-critical robotic control.
        \item \textbf{Multi-Modal Visual Question Answering}: Connecting camera image features directly into the on-device language model to enable interactive
        visual troubleshooting for machine operators.
    \end{enumerate}
    
    \section*{Acknowledgments}
    The authors express their sincere gratitude to the Department of Computer Science \& Engineering at MIT-ADT University, Pune, and faculty guide Prof. Dr.
    Harshad Lokhande, for their continuous support, laboratory facilities, and technical guidance throughout this research.
    
    \begin{thebibliography}{00}
    \bibitem{b1} W. Shi, J. Cao, Q. Zhang, Y. Li, and L. Xu, ``Edge Computing: Vision and Challenges,'' \textit{IEEE Internet of Things Journal}, vol. 3, no. 5,
    pp. 637--646, 2016.
    \bibitem{b2} F. Bonomi, R. Milito, J. Zhu, and S. Addepalli, ``Fog Computing and Its Role in the Internet of Things,'' \textit{Proc. 1st Edition of the MCC
    Workshop on Mobile Cloud Computing}, pp. 13--16, 2012.
    \bibitem{b3} OpenFog Consortium, ``OpenFog Reference Architecture for Fog Computing,'' 2017. (Adopted as IEEE Standard 1934-2018).
    \bibitem{b4} M. Satyanarayanan, ``The Emergence of Edge Computing,'' \textit{Computer}, vol. 50, no. 1, pp. 30--39, 2017.
    \bibitem{b5} B. Varghese and R. Buyya, ``Next Generation Cloud Computing: New Trends and Research Directions,'' \textit{Future Generation Computer Systems},
    vol. 79, pp. 849--861, 2018.
    \bibitem{b6} B. Burns, B. Grant, D. Oppenheimer, E. Brewer, and J. Wilkes, ``Borg, Omega, and Kubernetes,'' \textit{Communications of the ACM}, vol. 59, no. 5,
    pp. 50--57, 2016.
    \bibitem{b7} D. Merkel, ``Docker: Lightweight Linux Containers for Consistent Development and Deployment,'' \textit{Linux Journal}, vol. 2014, no. 239, 2014.
    \bibitem{b8} C. Pahl and P. Jamshidi, ``Microservices: A Systematic Mapping Study,'' \textit{Proc. 6th Int. Conf. on Cloud Computing and Services Science}, pp.
    137--146, 2016.
    \bibitem{b9} Amazon Web Services, ``AWS IoT Greengrass Architecture and Developer Guide,'' AWS Whitepapers, 2024.
    \bibitem{b10} Microsoft, ``Azure IoT Edge Architecture and Security Documentation,'' Microsoft Docs, 2024.
    \bibitem{b11} D. Thain, T. Tannenbaum, and M. Livny, ``Distributed Computing in Practice: The Condor Experience,'' \textit{Concurrency and Computation:
    Practice and Experience}, vol. 17, no. 2--4, pp. 323--356, 2005.
    \bibitem{b12} WebAssembly Community Group, ``WebAssembly Core Specification and Component Model,'' W3C Recommendation, 2024.
    \bibitem{b13} A. K. Paul et al., ``Comparative Analysis of Lightweight Kubernetes Distributions for Edge Computing: Performance and Resource Efficiency,''
    \textit{arXiv preprint arXiv:2504.03656}, 2025.
    \bibitem{b14} W. Kwon et al., ``Efficient Memory Management for Large Language Model Serving with PagedAttention,'' \textit{Proc. 29th ACM Symp. on Operating
    Systems Principles (SOSP)}, 2023.
    \bibitem{b15} G. Gerganov et al., ``llama.cpp: Fast LLM Inference in C/C++ without Dependencies,'' GitHub Repository, 2024.
    \bibitem{b16} J. Redmon, S. Divvala, R. Girshick, and A. Farhadi, ``You Only Look Once: Unified, Real-Time Object Detection,'' \textit{Proc. IEEE Conf.
    Computer Vision and Pattern Recognition (CVPR)}, pp. 779--788, 2016.
    \bibitem{b17} G. Jocher, A. Chaurasia, and J. Qiu, ``Ultralytics YOLOv8 and YOLO11 Architecture Documentation,'' GitHub Repository, 2024.
    \bibitem{b18} Microsoft, ``ONNX Runtime: Cross-Platform High Performance ML Inferencing Accelerator,'' ONNX Runtime Documentation, 2024.
    \bibitem{b19} International Organization for Standardization, ``Mechanical Vibration---Evaluation of Machine Vibration by Measurements on Non-Rotating Parts,''
    \textit{ISO Standard 10816-1}, 2016.
    \bibitem{b20} O. Abdeljaber et al., ``Real-Time Vibration-Based Structural Damage Detection Using One-Dimensional Convolutional Neural Networks,''
    \textit{Journal of Sound and Vibration}, vol. 388, pp. 154--170, 2017.
    \bibitem{b21} Y. Koizumi et al., ``Acoustic Anomaly Detection in Machines Using Mel-Frequency Cepstral Coefficients and Autoencoders,'' \textit{Proc. IEEE Int.
    Conf. Acoustics, Speech and Signal Processing (ICASSP)}, pp. 855--859, 2019.
    \bibitem{b22} M. Manso, J. V. Silva, and R. Santos, ``Secure and Reliable Over-The-Air Firmware Updates for Industrial Internet of Things,'' \textit{IEEE
    Internet of Things Magazine}, vol. 4, no. 3, pp. 62--67, 2021.
    \end{thebibliography}

    \end{document}